\subsection{The exponent-nineteen quotient and its arithmetic criterion}
\label{subsec:nineteen-arithmetic}

Let $C$ be the smooth projective curve with affine equation
\[
 y^{19}=x(1-x)^7,
\]
and put $J=\operatorname{Jac}(C)$.  Write
\[
 K=\mathbb Q(\zeta_{19}),\quad H=\langle7\rangle
 \subset\operatorname{Gal}(K/\mathbb Q),\quad
 K_0=K^H,\quad F=K^+,\quad F_0=K_0^+.
\]
Thus $[K:\mathbb Q]=18$, $[K_0:\mathbb Q]=6$,
$[F:\mathbb Q]=9$, and $[F_0:\mathbb Q]=3$.

\begin{proposition}[A sufficient simple-zero certificate]
\label{prop:lm-nineteen-rankone}
There are abelian varieties $A,P$ over $\mathbb Q$ of dimensions
three and six respectively, and a rational isogeny
$J\sim_{\mathbb Q} A\times P$, such that
\[
 \operatorname{rank} A(\mathbb Q)=1,
 \qquad \operatorname{rank} P(\mathbb Q)=0.
\]
If
\[
 L'(J/\mathbb Q,1)\ne0,
\]
then $\Sha(J/\mathbb Q)$ and $\Sha(P/\mathbb Q)$ are finite.
Consequently this analytic certificate supplies a nonconstant map
from the Fermat curve of exponent nineteen to a rank-zero abelian
variety with finite Tate--Shafarevich group.
\end{proposition}

\begin{proof}
The rational automorphism
\[
 c(x,y)=\left(\frac1{1-x},-\frac{y^7}{(1-x)^3}\right)
\]
has order three.  Indeed its second iterate is
\[
 c^2(x,y)=\left(\frac{x-1}{x},
                 -\frac{y^{11}}{x(1-x)^4}\right),
\]
and its third iterate is the identity, using $y^{19}=x(1-x)^7$.
The map on the $x$-line has the two fixed points
$x^2-x+1=0$.  Above either fixed point its action on the
nineteen-element geometric fiber has multiplier $7$ in
$\mathbb F_{19}$, hence exactly one fixed point.  It cyclically
permutes the three branch points, so there are no further fixed
points.  Riemann--Hurwitz gives genus three for
$B=C/\langle c\rangle$.  Put $A=\operatorname{Jac}(B)$ and
\[
 e=\frac{1+c+c^2}{3}\in\operatorname{End}^0_{\mathbb Q}(J).
\]
The image of $e$ is isogenous to $A$, while the image of $1-e$
defines the six-dimensional complementary factor $P$.

The curve $C$ over $\mathbb Q$ is isomorphic to the curve
$F_{7,1,-8}$ in \citet[Theorem~1.6(2), author p.~3]{McCallumTzermias2004},
by interchanging $x$ and $1-x$. Their theorem proves
$\operatorname{rank}J(\mathbb Q)=1$.
On the one-dimensional space $J(\mathbb Q)\otimes\mathbb Q$,
the order-three operator $c$ must act trivially, since
$\mathbb Q^\times$ has no nonidentity element of order three.
Thus the entire rank is in the image of $e$, proving the two
rank assertions.  This does not yet assert finiteness of any
full Tate--Shafarevich group.

We identify the Hilbert form needed for that finiteness assertion.
Let $\tau(x,y)=(x,\zeta_{19}y)$.  The relation
$c\tau c^{-1}=\tau^7$ gives an action over $K$ of the crossed
product algebra
\[
 K\rtimes H\simeq\operatorname{End}_{K_0}(K)
                  \simeq M_3(K_0)
\]
on $J$.  This algebra is simple, and its homomorphism to
$\operatorname{End}^0_K(J)$ is nonzero, so it is injective.
Here the generator of $H$ acts as $c$, with $c^3=1$; thus the
crossed product has the trivial factor set.
Complex conjugation on the field of definition takes
$\tau$ to $\tau^{-1}$ and fixes $c$.  Under the displayed
identification this is the conjugation on endomorphisms induced
by the involution on $K$.  Since $K=FK_0$, taking its fixed
subalgebra gives
\[
 \operatorname{End}_{K_0}(K)^{\operatorname{Gal}(K/F)}
       =\operatorname{End}_{F_0}(F)\simeq M_3(F_0)
       \subset\operatorname{End}^0_F(J).
\]
The idempotent $e$ is the rank-one trace projector in this
matrix algebra.  Matrix units therefore give an $F$-isogeny
\[
 J_F\sim_F A_F^3.
\]
This is an isogeny over the degree-nine real field, not an
unproved rational splitting into three copies.

For completeness, the CM types and Hecke coefficients also
identify the simple factor.  The CM type of the Jacobi-sum
character is
\[
 \{1,2,3,4,6,7,9,11,14\}\subset(\mathbb Z/19\mathbb Z)^\times.
\]
Its stabilizer is exactly $H$, by \citet[Lemma~1.6, pp.~206--207]{GrossRohrlich1978}.  Consequently the induced type on
$K_0$ is primitive and $A$ is absolutely simple with CM by
$K_0$.  The Jacobi sum indexed by $(1,7,11)$ is unchanged
when all three indices are multiplied by $7$, since this
cyclically permutes the triple.  Its values therefore lie in
$K_0$.  The real quadratic-induction traces lie in $F_0$.

More explicitly, after extending the $\ell$-adic coefficient field,
write $\psi_a$ for the CM character indexed by
$a\in(\mathbb Z/19\mathbb Z)^\times$ on $G_K$.
The Jacobi-sum symmetry gives $\psi_{7a}=\psi_a$.
The three representations
\[
 W_a=\operatorname{Ind}_{G_K}^{G_F}\psi_a,
 \qquad a\in(\mathbb Z/19\mathbb Z)^\times/(\{\pm1\}H),
\]
are pairwise nonisomorphic: their restrictions to $G_K$ are
$\psi_a\oplus\psi_{-a}$, and the stabilizer of the CM type
is exactly $H$.  Each occurs three times in $V_\ell(J_F)$.
On these multiplicity spaces the operator $c$ is the regular
order-three action, so $e$ selects one copy of each.  Therefore
\[
 V_\ell(A_F)\otimes\overline{\mathbb Q}_\ell
       =\bigoplus_a W_a.
\]
This describes the three real-multiplication components explicitly.

Let $f$ be the Hilbert form over $F$ obtained by quadratic
automorphic induction of the Jacobi-sum Hecke character over
$K$.  The character is not a norm character, because its
two infinity exponents at each complex place differ by one.
Thus it gives a cuspidal parallel-weight-two form.  In unitary
normalization its restriction to $\mathbb A_F^\times$ is
$\eta_{K/F}$, so the determinant formula for induction makes
its finite central character trivial.  These are the explicit
CM and induction constructions in \citet[\S\S2.1--2.2 and 3.4, author pp.~5--6 and 15]{ShuExpanded2021}; they can equally be obtained
from \citet[Theorems~10--11]{SerreTate1968} and
\citet[Proposition~12.1]{JacquetLanglands1970}.

The standard Jacobi-sum Euler factors give
\[
 L(f/F,s)=L(J/\mathbb Q,s).
\]
Every coefficient conjugate of $f$ has the same $L$-series:
conjugating its Hecke character is the same as reindexing
the primes of the abelian field $K$, which preserves norms.
The Hecke field is $F_0$.  Indeed it is contained in $F_0$
by the trace calculation, and a proper subfield would be
$\mathbb Q$.  In that case the existence statement immediately
before Zhang's Theorem~A supplies an attached elliptic curve
$E/F$, since $[F:\mathbb Q]$ is odd.  All three coefficient
conjugates would have the same two-dimensional system
$V_\ell(E)$, so the preceding representation identity and
the isogeny criterion of \citet[\S5, Corollary~2, p.~22]{Faltings1986} would give $A_F\sim_F E^3$,
contradicting the absolute simplicity of $A$.
It follows that $A_F$ is the abelian variety attached to
$f$, with
\[
 L(A/F,s)=\prod_{\iota:F_0\hookrightarrow\mathbb R}L(f^\iota/F,s)
          =L(J/\mathbb Q,s)^3.
\]
The equality here is compatible with the preceding isogeny;
in particular one does not apply a rank-one theorem to the
degree-fifty-four $L$-function of $A/F$ itself.

\citet[Theorem~A, p.~28; proof in \S7.2, p.~143]{Zhang2001}
applies to a parallel-weight-two Hilbert newform with trivial
central character if the totally real degree is odd or a
finite local conductor has exponent one.  It states that
analytic order at most one gives finite full
Tate--Shafarevich group for the associated abelian variety.
There is no CM exclusion in that theorem.  Here $[F:\mathbb Q]=9$
is odd, so no auxiliary finite-place condition is needed.
The root number of $L(J/\mathbb Q,s)$ is $-1$, as verified
below.  The assumed nonzero derivative therefore makes
$f$ have analytic order exactly one.  Zhang's theorem gives
$\Sha(A/F)$ finite.

Finiteness of the full Tate--Shafarevich group is invariant
under isogeny over a number field.  Hence the $F$-isogeny
$J_F\sim A_F^3$ gives $\Sha(J/F)$ finite.
The kernel of restriction
$\Sha(J/\mathbb Q)\to\Sha(J/F)$ lies in
\[
 H^1(\operatorname{Gal}(F/\mathbb Q),J(F)).
\]
This group is finite: $J(F)$ is finitely generated, the
Galois group is finite, and its positive-degree cohomology
is killed by the group order.  Thus $\Sha(J/\mathbb Q)$ is
finite.  The rational isogeny $J\sim A\times P$ gives
$\Sha(P/\mathbb Q)$ finite as well.

Finally, the quotient map from the exponent-nineteen Fermat
curve to $C$, followed by a rational-cusp Abel map and a
nonzero integral multiple of $1-e$, is nonconstant into $P$.
Indeed the Abel image of $C$ generates $J$, so its image
under a surjection onto the positive-dimensional factor
$P$ generates $P$ and cannot be constant.
\end{proof}

\paragraph{Exact analytic input to be certified.}
For the triple $(r,s,t)=(1,7,11)$, the residue
\[
 r^r s^s(t-19)^t\equiv68\pmod {19^2}
\]
has Fermat quotient zero modulo $19$.  Equivalently the
sum of inverse representatives of the displayed CM type
is $95\equiv0\pmod {19}$.  \citet[Proposition~3.12(1), author p.~16]{ShuExpanded2021} gives
Hecke conductor $\mathfrak P$, and therefore the induced
rational conductor is
\[
 N=|\operatorname{disc}K|\,N_{K/\mathbb Q}\mathfrak P
   =19^{17}\cdot19=19^{18}.
\]
\citet[Theorem~3.13, author pp.~17--18]{ShuExpanded2021} gives the tame global sign
$\left(\frac2{19}\right)=-1$.
The completed function is
\[
 \Lambda(s)=N^{s/2}(2\pi)^{-9s}\Gamma(s)^9L(J/\mathbb Q,s),
 \qquad \Lambda(s)=-\Lambda(2-s).
\]
Thus a rigorous interval excluding zero for
$L'(J/\mathbb Q,1)$, or equivalently for $\Lambda'(1)$,
is sufficient. Proposition~\ref{prop:nineteen-analytic-certificate}
proves such an interval using three degree-six factors.

\paragraph{The required character order.}
The subgroup $H$ of order three lies in the kernel of every
cubic Dirichlet character modulo nineteen, since that kernel
has order six.  Characters of order nine do restrict nontrivially to $H$
and give the factorization used in the next subsection. The arithmetic
criterion itself used the crossed-product algebra and odd-degree
real base change, independently of that analytic factorization.

